\chapter{Discretionary Macro}\label{ch:s2:discretionary-macro}

A macro manager with a view that rates will fall has a dozen ways to express it. The view is the easy part; the instrument, the size and the stop
decide whether being right pays. This chapter holds one view fixed and changes everything else. The manager expects a rate to fall 50 basis points
over six months, and on the synthetic paths the rate ends lower 73.4\% of the time. Expressed through futures with a stop 25 basis points away, the
view is stopped out on 26.5\% of the paths where it turns out right. With a stop 50 basis points away that falls to 5.6\%. Expressed through an
at-the-money option, 17.2\% of correct views still lose the premium. The build is \texttt{firm.macrobook}. The chapter has no strategy files: its
subject is how any macro view becomes a position.

\section{From view to trade}

\begin{definition}[Trade expression]\label{def:s2:discretionary-macro:expression}
A \emph{trade expression}\index{trade expression} is the instrument, strike, maturity and size chosen to turn a market view into a position; different
expressions of the same view pay differently depending on how far, how fast and along which path the market moves.
\end{definition}

A view that ``rates will fall'' has a size (how much), a horizon (by when), a path (steadily or after a rise) and a confidence. Each expression bets
on a different combination. A futures position bets on the size, whatever the path. An option bets on the size and caps the loss at its premium. A
futures position with a stop bets that the path does not first go the wrong way. A spread bets on a move up to a point and gives up the rest.

\texttt{firm.macrobook} makes this concrete (\cref{lst:s2:discretionary-macro:expressions}). The manager's view is a 50 basis-point fall over six
months. The rate's true drift is drawn around the view with a dispersion of 50 basis points, so the view is right on average but by a varying amount.
The rate moves with a normal volatility of 90 basis points a year. Over 20,000 paths it ends lower on 73.4\% of them.

\section{Instrument choice and convexity}

\begin{definition}[Convexity budget]\label{def:s2:discretionary-macro:convexity}
A \emph{convexity budget}\index{convexity budget} is the premium a manager is willing to spend on options to express views, whose most it can lose;
spending it buys positions whose loss is known in advance and whose gain grows with the size of the move.
\end{definition}

Each expression is sized to the same risk budget. For options the budget is the premium. For futures with a stop it is the stop distance. For futures
without a stop it is a two-standard-deviation loss over the horizon. Options are priced by the normal model at the rate's volatility plus 10\%. An
at-the-money receiver option costs 27.9 basis points, one struck 25 basis points below costs 17.2, and an option spread between the two strikes 0 and
$-50$ costs 18.2 (\cref{fig:s2:discretionary-macro:expressions}).

\begin{center}
\small
\begin{tabular}{@{}lccccc@{}}
\toprule
20,000 paths, per unit of risk budget & mean & sd & P(profit) & P(loss $|$ right) & 5th percentile\\
\midrule
futures, no stop & 0.39 & 0.63 & 73\% & 0.0\% & $-0.65$\\
futures, stop 25 bp & 1.68 & 2.98 & 54\% & 26.5\% & $-1.00$\\
futures, stop 50 bp & 0.98 & 1.55 & 69\% & 5.6\% & $-1.00$\\
option at the money & 1.26 & 2.23 & 61\% & 17.2\% & $-1.00$\\
option 25 bp out of the money & 1.69 & 3.20 & 54\% & 26.5\% & $-1.00$\\
option spread (0 to $-50$) & 0.71 & 1.22 & 65\% & 10.8\% & $-1.00$\\
\bottomrule
\end{tabular}
\end{center}

The ranking depends on what the budget is. Per unit of a two-standard-deviation loss, futures earn 0.39 and never lose when the view is right. Per unit
of premium, the out-of-the-money option earns 1.69. That figure is mostly leverage, with the same 26.5\% of correct views ending worthless as the tight
stop. The option spread is the most conservative option: it gives up the large moves, which the view expects less of, and wins more often (65\%). No
expression is better than another in general. Each is right for a different belief about how the view will come true.

\begin{omfigure}
\begin{tikzpicture}
\begin{axis}[width=12cm,height=5.6cm,ybar,bar width=11pt,ylabel={\small per unit of budget / share},tick label style={font=\footnotesize},
  xmin=0.4,xmax=6.6,ymin=0,ymax=1.9,xtick={1,2,3,4,5,6},
  xticklabels={futures,{stop 25},{stop 50},{option ATM},{option OTM},{spread}},x tick label style={font=\scriptsize,rotate=20,anchor=north east},
  grid=major,grid style={gray!25},table/col sep=comma,legend style={font=\scriptsize,at={(0.5,-0.36)},anchor=north},legend columns=2]
\addplot[fill=omThm!70,draw=omThm] table[x=x,y=mean]{figdata/strategies-2/17-discretionary-macro/expressions.csv};
\addplot[fill=gray!40,draw=gray] table[x=x,y=lose]{figdata/strategies-2/17-discretionary-macro/expressions.csv};
\legend{mean payoff per unit of risk budget,share of correct views that lose}
\end{axis}
\end{tikzpicture}

\omcaption{One view (a 50 basis-point fall in six months) expressed six ways on 20,000 synthetic paths: the mean payoff per unit of each expression's
risk budget, and the share of paths on which the view was right but the expression lost. Data: \texttt{s2\_macrobook.table}.}\label{fig:s2:discretionary-macro:expressions}
\end{omfigure}

\section{Sizing and stops}

\begin{definition}[Stop discipline]\label{def:s2:discretionary-macro:stop}
\emph{Stop discipline}\index{stop discipline} is the rule, set when a position is opened, for cutting it after a given loss or adverse move, and the
practice of following the rule; it limits the loss on wrong views at the price of cutting some right ones before they pay.
\end{definition}

A stop turns the path into the risk. The rate can rise 25 basis points on the way to falling 50, and a stop at 25 closes the position before the fall.
The share of correct views stopped out falls steeply with the stop's distance (\cref{fig:s2:discretionary-macro:stops}): 55.6\% at 10 basis points,
26.5\% at 25, 5.6\% at 50, 0.7\% at 75 and 0.1\% at 100. Kaminski and Lo gave a framework for when stop-loss rules add or subtract expected return
and volatility. In their empirical analysis of index futures, some stop-loss policies raised expected return while cutting volatility at longer
sampling frequencies. In the synthetic market the drift is the view, so the stop pays only when the view is wrong.

\begin{omfigure}
\begin{tikzpicture}
\begin{axis}[width=10cm,height=5.4cm,xlabel={\small stop distance (bp)},ylabel={\small correct views stopped out (\%)},tick label style={font=\footnotesize},
  xmin=0,xmax=100,ymin=0,ymax=75,grid=major,grid style={gray!25},table/col sep=comma]
\addplot[omThm,thick,mark=*,mark size=1.3pt] table[x=stop,y=share]{figdata/strategies-2/17-discretionary-macro/stops.csv};
\end{axis}
\end{tikzpicture}

\omcaption{The share of correct views (the rate ends lower after six months) whose futures position was stopped out first, by the stop's distance,
on the synthetic paths. Data: \texttt{s2\_macrobook.stops}.}\label{fig:s2:discretionary-macro:stops}
\end{omfigure}

Sizing follows from the budget and the stop. With the same budget, a tighter stop means a larger position: the stop at 25 basis points holds twice
the position of the stop at 50. That is why its mean payoff per unit of budget is higher (1.68 against 0.98) while it is stopped out far more often. A
manager whose budget is a monthly loss limit (Book~8, chapter~28, on drawdown limits) must choose between a small position with a wide stop and a
large one with a tight stop. The first survives noise; the second is right less often but pays more when it is.

\section{Keeping score}

A discretionary book is judged on outcomes, but its decisions are made on views. Separating the two is the point of keeping score. Record each view
when it is formed: size, horizon, confidence. Record the expression chosen and the reason. When the trade closes, record whether the view came true,
whether the expression paid, and whether the stop or the horizon ended it.

Over many trades the record shows whether the manager's views are better than chance, and whether the expressions make the most of them. In the
synthetic market with the view known exactly (a dispersion of zero), the rate ends lower on 78.6\% of the paths: the rest is noise within six months.
Even a perfectly informed view is right only that often at this horizon. The tight-stop futures then earn 1.45 per unit of budget and the at-the-money
option 1.07. A score of right and wrong alone would miss that.

\section{Tutorial: right and still losing}\label{tut:s2:discretionary-macro:tut}

\begin{tutorial}
\textbf{Goal.} Simulate a view's paths, express it six ways, and measure how often a correct view loses and how stops cut it short.
\textbf{End state:} the table and the two figures.
\begin{tutsteps}
\item \textbf{Expressions}.
\omcode{code/firm/macrobook/firm_macrobook.py}{60}{82}{Each expression's payoff per unit of risk budget, and correct views stopped out.}\label{lst:s2:discretionary-macro:expressions}
\item \textbf{The table}.
\omcode{code/strategies-2/17-discretionary-macro/python/s2_macrobook.py}{30}{39}{Mean, spread, chance of profit and of losing when right.}\label{lst:s2:discretionary-macro:table}
\item \textbf{Run} \texttt{table()}, \texttt{stops()}, \texttt{premiums()} and \texttt{fig\_macrobook.py}.
\end{tutsteps}
\textbf{What to change next.} Make the view's timing uncertain (the fall comes in the second half); add a trailing stop; express the view as a
curve trade in a two-factor model.
\end{tutorial}

\section{Build: macro book}\label{bld:s2:discretionary-macro:macrobook}

\begin{build}
\bfield{Purpose} A view's distribution of paths, expressions sized to a common risk budget, and stop rules with their cost.
\bfield{Interface} \texttt{ViewConfig(\dots)}, \texttt{rate\_paths(cfg)}, \texttt{bachelier\_receiver(strike, vol, tau)},
\texttt{expressions(paths, cfg, stops)}, \texttt{stopped\_correct(paths, cfg, stop)}.
\bfield{Rules} Options priced at the path's volatility times a premium; each expression's payoff divided by its own risk budget.
\bfield{Acceptance tests} \texttt{code/firm/macrobook/tests/}: the normal-model price at the money and put--call parity; expressions and stops on
hand-made paths.
\bfield{Stretch} Uncertain timing; trailing stops; two-factor curves.
\end{build}

\begin{omsources}
\item K.~M.~Kaminski and A.~W.~Lo, ``When do stop-loss rules stop losses?'', \emph{Journal of Financial Markets} 18, 2014.
\end{omsources}

\section{Exercises}

\begin{exercise}[$\star$]\label{exo:s2:discretionary-macro:1}
A receiver option at the money on a rate with normal volatility 99 basis points a year has six months to run. What is its price?
\end{exercise}

\begin{exercise}[$\star$]\label{exo:s2:discretionary-macro:2}
A budget of one unit buys futures with a 25 basis-point stop or with a 50 basis-point stop. How do the two positions compare in size?
\end{exercise}

\begin{exercise}[$\star$]\label{exo:s2:discretionary-macro:3}
The rate falls 80 basis points. What does the option spread (0 to $-50$) pay per unit of its 18.2 basis-point premium?
\end{exercise}

\begin{exercise}[$\star\star$]\label{exo:s2:discretionary-macro:4}
Why is the share of correct views stopped out so sensitive to the stop's distance?
\end{exercise}

\begin{exercise}[$\star\star$]\label{exo:s2:discretionary-macro:5}
Why does the out-of-the-money option have a higher mean payoff per unit of budget and a lower chance of profit than the at-the-money one?
\end{exercise}

\begin{exercise}[$\star\star$]\label{exo:s2:discretionary-macro:6}
Why is a view that is exactly right still wrong on 21\% of paths at six months?
\end{exercise}

\begin{exercise}[$\star\star\star$]\label{exo:s2:discretionary-macro:7}
\emph{Coding.} Run \texttt{table(0.0)}, a view with no dispersion. How do the chances and payoffs change, and what does that say about the value of
precision?
\end{exercise}

\begin{exercise}[$\star\star\star$]\label{exo:s2:discretionary-macro:8}
\emph{Find the flaw.} ``I was right about rates this year and still lost money, so my views are fine and the market was wrong.''
\end{exercise}

\section{Problem: Right and Still Losing}

\begin{problem}[{Weekend problem --- expressing a macro view}]\label{pb:s2:discretionary-macro:1}
The chapter's synthetic paths and the public record.

\medskip\noindent\textbf{Part I --- The view.}
\begin{enumerate}
\item Define a \omterm{def:s2:discretionary-macro:expression}{trade expression}.
\item What does each expression bet on?
\item Describe the synthetic view and its paths.
\item How often is the view right?
\end{enumerate}

\medskip\noindent\textbf{Part II --- Instruments.}
\begin{enumerate}[resume]
\item Define a \omterm{def:s2:discretionary-macro:convexity}{convexity budget}.
\item How is each expression sized?
\item Give the options' prices.
\item Why does the ranking depend on the budget?
\end{enumerate}

\medskip\noindent\textbf{Part III --- Stops.}
\begin{enumerate}[resume]
\item Define \omterm{def:s2:discretionary-macro:stop}{stop discipline}.
\item Give the share of correct views stopped out by stop distance.
\item What did Kaminski and Lo study and find?
\item How do stop and size interact?
\end{enumerate}

\medskip\noindent\textbf{Part IV --- The verdict.}
\begin{enumerate}[resume]
\item State the \emph{named result}: the expected payoff of each expression of the same view and the share of correct views stopped out early.
\item Which expression suits a view about size, and which one about path?
\item What changes when the view is exact?
\item What should a manager record for each trade?
\item How would you judge a discretionary manager?
\item Why does this chapter have no strategy files?
\item How does this chapter relate to chapter~16?
\item In one sentence: what decides whether a right view pays?
\end{enumerate}
\end{problem}

\section{Interview questions}

\begin{interviewq}[$\star$ \iqroles{trader}]\label{iq:s2:discretionary-macro:1}
You think the Fed will cut rates in six months. How would you express it?
\end{interviewq}

\begin{interviewq}[$\star\star$ \iqroles{trader}]\label{iq:s2:discretionary-macro:2}
When would you use options rather than futures for a macro view?
\end{interviewq}

\begin{interviewq}[$\star\star$ \iqroles{risk}]\label{iq:s2:discretionary-macro:3}
How would you set stops for a discretionary macro book?
\end{interviewq}

\begin{interviewq}[$\star\star$ \iqroles{researcher}]\label{iq:s2:discretionary-macro:4}
How would you measure whether a manager's views have skill?
\end{interviewq}

\begin{interviewq}[$\star\star$ \iqroles{developer}]\label{iq:s2:discretionary-macro:5}
Design a trade journal that records views, expressions and outcomes.
\end{interviewq}

\begin{interviewq}[$\star\star\star$ \iqroles{researcher}]\label{iq:s2:discretionary-macro:6}
For a driftless Brownian motion with volatility $\sigma$, derive the probability of touching $+s$ before time $T$, and use it to explain the shape of
the chapter's stop-out curve.
\end{interviewq}
